 %---------------------------------------------------------------
% SLO: angleški povzetek
% ENG: english abstract
%---------------------------------------------------------------
\selectlanguage{english} % Preklopi na angleški jezik
\addcontentsline{toc}{chapter}{Abstract}
\chapter*{Abstract}

\noindent\textbf{Title:} \ttitleEn
\bigskip

Extracting information from structured data using large language models requires balancing answer quality,
response latency, and resource consumption. This thesis investigates how the extraction method and data
serialization format influence these requirements. We compared direct context inference, in which models
answer questions from serialized database contents, with text-to-SQL, in which models generate queries
executed by a database engine. The evaluation used \num{10} locally deployed open-weight models,
\num{6} serialization formats, and databases from Spider and a subsampled version of BIRD. We measured
reference-answer agreement, token consumption, and response latency. Text-to-SQL achieved higher mean
agreement than every direct-inference format, exceeding the strongest direct format in normalized exact
match by \num{46.5} percentage points on Spider and \num{44.7} percentage points on sampled BIRD.
In separate comparisons on shared context-fit workloads, it also required fewer input tokens and had
lower mean response latency. For direct inference on flat relational tables, CSV required the fewest
input tokens and had the lowest mean latency, with Token-Oriented Object Notation (TOON) a close
alternative; more verbose formats offered no consistent advantage in answer agreement. These findings
inform a decision framework that prioritizes extraction-method selection before serialization
optimization and supports text-to-SQL as the preferred method when SQL execution is available and
the task can be expressed as a query.

\subsection*{Keywords}
\textit{\tkeywordsEn}
\clearemptydoublepage
